\section{Experiments}
\label{sec:exp}
\subsection{Setup}
All runs use one laptop with an NVIDIA RTX 4060 Laptop GPU (8\,GB), PyTorch 2.11
and CUDA 12.8. The 241-class stage trains \cnPreEpochs{} epochs over
\cnPreTrainN{} images at 823\,s per epoch, 69 minutes in total; the 21-way head
trains \cnEpochs{} epochs over \cnTrainN{} images at 282\,s per epoch, 38 minutes
in total. One epoch of the second stage was lost to machine standby and is
excluded from the timing but not from the weights, because the checkpoint
selected by validation accuracy is unaffected. Every configuration uses the same
seed; the spread of a re-run with a different seed on the Western baseline is
0.5 points, which we use as the noise floor throughout. Thresholds and the
temperature are fitted on the validation split and then frozen.

\subsection{Pre-Training on Chinese Dishes}
\label{sec:pretrain}
Table \ref{tab:cnmain} reports the canteen head with and without the 241-class
stage. The two runs differ only in the initialisation of the trunk, so the
\cnGainPretrain-point difference is the value of Chinese pre-training under this
recipe: \cnNoPreTest\% against \cnTest\% on \cnTestN{} test images, with validation
accuracy moving from \cnNoPreVal\% to \cnVal\%. The gain is larger than any
adaptation-policy change we could obtain on the Western benchmark, where the same
kind of intervention (unfreezing one stage) is worth 6.4 points and the largest
data intervention is worth 1.2. It is also the cheapest: the pre-training stage
costs 69 minutes once and adds no parameters and no inference cost.

What pre-training does not fix is the boundary of the catch-all class. The
twenty-dish rows of the confusion matrix reach \cnDishAcc\% while the catch-all
bucket is recalled at \cnOtherRecall\%, and the gap between those two numbers is
the honest measure of how far a menu-only classifier is from a deployed one.

\begin{table}[t]
\centering
\caption{The 21-way canteen head. Val is the split used for checkpoint
selection; test is the held-out split of \cnTestN{} images. ECE is computed with
fifteen bins on the single-view protocol.}
\label{tab:cnmain}
\small
\begin{tabular*}{\columnwidth}{@{}lrrrr@{}}
\toprule
Configuration & Val & Test & Top-5 & ECE \\
\midrule
No Chinese pre-training & \cnNoPreVal & \cnNoPreTest & \cnNoPreTopFive & \cnNoPreEce \\
Chinese pre-training (deployed) & \cnVal & \cnTest & \cnTopFive & \cnEce \\
\quad + stronger photometric aug. & \ablStrongVal & \ablStrongTest & \ablStrongTopFive & \ablStrongEce \\
\quad + class-balanced loss & \ablBalVal & \ablBalTest & \ablBalTopFive & \ablBalEce \\
\quad + balanced, catch-all unweighted & \ablBalNoOtherVal & \ablBalNoOtherTest & \ablBalNoOtherTopFive & \ablBalNoOtherEce \\
\quad + focal loss ($\gamma=1$) & \ablFocalVal & \ablFocalTest & \ablFocalTopFive & \ablFocalEce \\
\bottomrule
\end{tabular*}
\end{table}

\subsection{Six-View Fusion and Per-Class Error}
Table \ref{tab:cnfuse} separates the two protocols. The single-view number is
what checkpoint selection optimises and what the training script reports; the
six-view number is what the service actually computes. Fusion is worth
\cnFuseGain{} points, which is more than the seed spread, and it lifts top-5 to
\cnFuseTopFive\%: at that point the correct dish is in the top three for almost
every photograph that the gate accepts.

Figure \ref{fig:cnconf} shows where the remaining error sits. The diagonal is
clean for the dishes with a distinctive colour or texture---potato shreds,
di san xian, tomato and egg, steamed egg all above 94\%---and it degrades exactly
where the menu is redundant: farmhouse pork against pepper pork (75.5\% and
77.1\% recall, with fifteen farmhouse images answered as pepper pork),
sweet-and-sour ribs against sweet-and-sour pork, braised chicken against braised
pork. The catch-all row is the most informative: its mass leaks into farmhouse
pork, fried noodles, stir-fried greens and braised chicken, i.e.\ into the
generic-looking stir-fries, which is what one expects when 175 categories are
compressed into one bucket.

\begin{table}[t]
\centering
\caption{Single-view protocol (used for model selection) against the six-view
fusion used by the service, on \cnTestN{} test images. The last three columns
are computed from the fused confusion matrix.}
\label{tab:cnfuse}
\small
\begin{tabular*}{\columnwidth}{@{}lrrrrr@{}}
\toprule
Protocol & Top-1 & Top-5 & 20 dishes & Catch-all & Macro-F1 \\
\midrule
Single view & \cnTest & \cnTopFive & -- & -- & -- \\
Six views & \cnFuseTop & \cnFuseTopFive & \cnDishAcc & \cnOtherRecall & \cnMacroFone \\
\bottomrule
\end{tabular*}
\end{table}

\begin{figure}[t]
\centering
\includegraphics[width=\columnwidth]{cn21_confusion.png}
\caption{Row-normalised confusion matrix of the 21-way canteen head on its test
split, six-view protocol. The red frame marks the catch-all row; its mass leaks
into the generic stir-fries rather than spreading uniformly.}
\label{fig:cnconf}
\end{figure}

\subsection{Open-Set Behaviour}
\label{sec:openset}
Table \ref{tab:gate} and Figure \ref{fig:gate} report acceptance on six sets.
Without the domain probe the gate is already useful: it rejects 85\% of non-food
scenes, 77\% of fifteen unseen Chinese dishes and 82\% of real photographs of
sixteen Western dishes, while keeping \gateInOff\% of the supported test set. The
domain probe then removes a further 7.2 points of unseen Chinese dishes and 9.9
points of non-food scenes, at a cost of \gateInCost{} points of coverage on the
supported set and \gateValCost{} points on validation.

The trade is the one a deployed system wants: a rejected photograph costs the
user one retry, an accepted wrong dish corrupts the energy estimate silently.
The cost is not uniform across the menu, and that is the honest weakness of the
design: on the supported set the probe is confident for most dishes, but it
rejects the ones whose photographs look least like a plated canteen dish, which
is why acceptance is 84\% on the test split and only 31\% on the web images of
baozi (Section \ref{sec:web}).

\begin{table}[t]
\centering
\caption{Acceptance rate (\%) of the deployed gate, with and without the domain
probe, at threshold \gateThr. Higher is better only for the first two rows.}
\label{tab:gate}
\small
\begin{tabular*}{\columnwidth}{@{}lrrr@{}}
\toprule
Evaluation set & w/o probe & deployed & $\Delta$ \\
\midrule
Canteen validation (\cnValN) & \gateValOff & \gateVal & $-$\gateValCost \\
Canteen test, 20 dishes (\cnDishN) & \gateInOff & \gateIn & $-$\gateInCost \\
Fifteen unseen Chinese dishes & \gateOutCnOff & \gateOutCn & $-$\gateOutCnGain \\
Non-Chinese food & \gateOutOodOff & \gateOutOod & $-4.5$ \\
Non-food and clutter & \gateOutOtherOff & \gateOutOther & $-9.9$ \\
Real photographs, 16 dishes & \gateRealOff & \gateReal & $-5.0$ \\
\bottomrule
\end{tabular*}
\end{table}

\begin{figure}[t]
\centering
\includegraphics[width=\columnwidth]{cn_gate.png}
\caption{What the system agrees to name, with and without the domain probe. The
supported sets stay above 90\%, the unsupported ones drop to 5--24\%. The
dashed line is 50\% and has no meaning of its own; it only makes the two groups
visually separable.}
\label{fig:gate}
\end{figure}

\subsection{External Evaluation on Web Images}
\label{sec:web}
The \webN{} web images are the only photographs in this paper that no part of
the pipeline has seen. The system is \webTop\% top-1 accurate on them and accepts
\webAccept\%, with \webTopAcc\% top-1 among the accepted ones. Two classes drag
the accuracy down: braised chicken at 75\% and braised pork at 84.6\%, both
confused with sweet-and-sour ribs, which is the same confusion the test split
shows and therefore a property of the classes rather than of the domain.

The acceptance rate is where the web set earns its place. It is \webAccept\%
against \gateIn\% on the test split overall, and the distribution is inverted:
baozi is
accepted 31.6\% of the time and poached chicken 61.5\%, because the domain probe
was fitted on photographs of plated dishes and web images of steamed buns and
poached chicken are shallow-depth-of-field studio shots with no table, no tray
and no context. Measured directly, the web images have a plain border in 0.3\%
of cases against 0.5\% for the test split but a median Laplacian variance of 802
against 505, i.e.\ they are sharper and more isolated than the photographs the
probe was trained on. Removing the domain probe would raise acceptance to 90.2\%
on this set, but it also raises non-food acceptance from 5.2\% to 15.1\%, so we
report the gap rather than tune the threshold on the very set we are reporting.
This is the one place where our external set is not representative: it measures
the menu, not the canteen.

\subsection{Negative Results}
The last four rows of Table \ref{tab:cnmain} are interventions that a reasonable
practitioner would try, and none of them survives.

\textbf{Stronger photometric augmentation} (crop scale 0.45 instead of 0.55,
jitter 0.5 instead of 0.35, hue 0.08, random grayscale with probability 0.1)
moves test accuracy from \cnTest\% to \ablStrongTest\% and top-5 to
\ablStrongTopFive\%, inside the seed spread, while costing three points on
poached chicken and four on mu xi pork. The failure mode is that canteen
variation is not photometric: the same dish is cut differently, garnished
differently and cooked to a different sauce colour, and none of those is a
colour jitter.

\textbf{Class-balanced loss} is the more interesting negative. Inverse-frequency
weights capped at four lift the twenty dishes from 87.4\% to 89.2\% and the macro
average from 86.3 to 88.0, and they fix precisely the weak classes: poached
chicken 82.4 to 94.1, sweet-and-sour ribs 71.0 to 77.9. But the catch-all bucket
loses 9.2 points, from 66.5\% to 57.2\%, and because it holds 36\% of the test
split the total falls to \ablBalTest\%. Pinning the catch-all weight back to one
restores it to 65.3\% and shrinks the dish gain to 0.6 points, netting
\ablBalNoOtherTest\%---a wash. The trade is therefore not in the weights but in
the decision boundary between "one of the twenty" and "something else", and
moving that boundary is exactly what the gate of Section \ref{sec:openset} is
for, with a knob that can be set per deployment instead of baked into the loss.

\textbf{Focal loss} ($\gamma=1$) trades accuracy for calibration: top-5 rises to
\ablFocalTopFive\% and ECE falls to \ablFocalEce, but validation accuracy drops
to \ablFocalVal\% and top-1 is unchanged. Since the displayed top-3 list and the
gate thresholds are both fitted on validation, we keep the cross-entropy model.

\subsection{Complexity, Latency and the Interface}
The deployed model is the ResNet-18 of the Western baseline, so its cost is
unchanged: \paramsFour\,M parameters, \macsFour\,GMACs at $224\times224$, and
\latGpuFour\,ms per forward pass on the GPU. Six views cost six forward passes,
which is still below 25\,ms on this GPU and around 100\,ms on the CPU. The
decision chain that surrounds the forward passes---gate, veto, catch-all test and
threshold comparison---takes \routeCn\,ms per request in the canteen tab against
\routeSingle\,ms in the single-dish tab, both measured on the same machine and
both dominated by the wide-head veto rather than by arithmetic.

The service exposes the three behaviours this paper measures. A supported
photograph returns the top three dishes with posteriors, the energy interval of
Eq.~\ref{eq:interval} and a portion selector; an unsupported one is refused with
the reason (not food, not a Chinese dish, another dish) and, when the refusal
comes from the domain or catch-all test rather than from a weak posterior, the
page offers to retry the same photograph in the single-dish or the 116-class tab.
That suggestion was evaluated offline on 19\,928 requests: when the canteen
tab refuses a real photograph the other tab names the right dish 94\% of the
time, and when it refuses a canteen test image the suggestion is right in only
10 of 300 cases, which is the expected asymmetry---the suggestion is for
photographs that were never canteen food.

\begin{figure}[t]
\centering
\includegraphics[width=0.95\columnwidth]{demo_cn.png}
\caption{The deployed interface on the canteen tab: top-3 dishes with posteriors,
energy intervals, portion control, and the gate's verdict in the response
payload.}
\label{fig:demo}
\end{figure}